%*******************************************************
% Abstract in English
%*******************************************************
\pdfbookmark[1]{Abstract}{Abstract}


\begin{otherlanguage}{american}
	\chapter*{Abstract}
	This thesis addresses the challenge of efficiently managing architectural information in Enterprise Architecture Management (EAM), where heterogeneous documentation and complex relationships require significant manual effort to extract and maintain. It introduces a Retrieval-Augmented Generation (RAG)-based chatbot named “\prototype{}”, designed to assist enterprise architects in extracting and understanding proprietary architectural data while maintaining transparency. \prototype{} connects to a knowledge graph containing both textbook-based EAM knowledge and user-uploaded enterprise architecture data. Natural language prompts are translated into database queries, and the retrieved information is exposed to the user along with the query used to augment the result. This increases the transparency of generated answers and reduces the black-box characteristics commonly associated with Large Language Model (LLM)-based systems.

Action Research is employed as the research methodology, enabling iterative development of the prototype in close collaboration between the researcher and two expert practitioners. Across four development cycles and a final qualitative review workshop, the system was evaluated regarding functionality, usability, answer quality, transparency, and alignment with practical EAM requirements.

The research demonstrates that transparent natural language interaction with a proprietary EAM knowledge graph is technically feasible. At the same time, the evaluation reveals limitations of the developed system regarding prompting robustness and operational maturity. This thesis contributes an implemented and empirically evaluated grey-box RAG architecture for EAM and highlights transparency, schema awareness, and practitioner involvement as decisive design principles when integrating LLMs with enterprise architectural data.
	
	\end{otherlanguage}